\section{Preliminaries}
\label{m:sec:prelim}

This section fixes notation and collects the standard results used later. Nothing
in it is new, and a reader fluent in continuous-time Markov chains and generating
functions may go straight to \cref{m:sec:gw}. It is included because four things are
used repeatedly and are easier to state once: that the competition between
exponential clocks converts rates into probabilities, that a system of forward
equations for state probabilities becomes a single partial differential equation
for a generating function, that conditioning on the first jump closes an equation
for hitting probabilities, and that the linear birth--death process is solvable by
exactly the route that \cref{m:sec:moc} will push further. Standard treatments are
Norris \cite{Norris1997}, Karlin and Taylor \cite{KarlinTaylor1975}, and, closest to
the applications here, Allen \cite{Allen2010}.

\subsection{Exponential clocks}
\label{m:sec:exp}

Every model in this thesis is built from events whose rate of occurrence does not
depend on how long the system has been waiting for them. That single assumption
forces the exponential distribution.

\begin{definition}[Exponential law]
\label{m:def:exp}
A \rv{} $T$ is exponentially distributed with rate $\theta>0$, written
$T\sim\mathrm{Exp}(\theta)$, if it has density and distribution function
\begin{equation*}
f_T(t) = \theta\ee^{-\theta t},
\qquad
F_T(t) = 1-\ee^{-\theta t},
\qquad t\ge0,
\end{equation*}
and $f_T(t)=0$ for $t<0$. Its mean is $1/\theta$ and its variance $1/\theta^2$.
\end{definition}

The property that matters is memorylessness. For $t_1,t_2>0$,
\begin{equation}
\label{m:eq:memoryless}
\pr{T > t_1+t_2 \mid T > t_1}
= \frac{\ee^{-\theta(t_1+t_2)}}{\ee^{-\theta t_1}}
= \ee^{-\theta t_2}
= \pr{T > t_2},
\end{equation}
so a clock that has already run for time $t_1$ is statistically indistinguishable
from a fresh one. The exponential law is the only continuous distribution with this
property, which is why a process assembled from constant-rate events depends on its
present state alone and not on its history.

When several events compete, the following two facts organise everything that
follows. Let $T_1,\dots,T_m$ be independent with $T_i\sim\mathrm{Exp}(\theta_i)$,
and let $T = \min_i T_i$ be the time of the first event.

\begin{proposition}[Competing clocks]
\label{m:prop:race}
With the notation above, write $\Theta = \sum_{i=1}^m\theta_i$. Then
\begin{equation}
\label{m:eq:race}
T \sim \mathrm{Exp}(\Theta),
\qquad
\pr{T_k = T} = \frac{\theta_k}{\Theta},
\end{equation}
and the identity of the winning clock is independent of $T$.
\end{proposition}

\begin{proof}
For the first claim,
$\pr{T>t} = \prod_i\pr{T_i>t} = \prod_i\ee^{-\theta_i t} = \ee^{-\Theta t}$.
For the second, condition on the value of $T_k$ and use independence:
\begin{equation*}
\pr{T_k = T}
= \int_0^\infty \theta_k\ee^{-\theta_k u}\prod_{i\neq k}\ee^{-\theta_i u}\,\dd u
= \theta_k\int_0^\infty \ee^{-\Theta u}\,\dd u
= \frac{\theta_k}{\Theta}.
\end{equation*}
Repeating the calculation with the integral truncated at $t$ gives
$\pr{T_k=T,\ T>t} = (\theta_k/\Theta)\ee^{-\Theta t}
= \pr{T_k=T}\pr{T>t}$, which is the asserted independence.
\end{proof}

\Cref{m:prop:race} is used in two ways. It says that a rate $\theta_k$ becomes a
probability $\theta_k/\Theta$ as soon as one asks which of several things happened
rather than when; \cref{m:sec:cts} uses exactly this to identify the replication
probability $p$ of a discrete-generation model with $\lambda/(\lambda+\mu)$ in a
continuous-time one. And it is the whole content of the direct simulation method.

\begin{remark}[Direct simulation]
\label{m:rem:gillespie}
\Cref{m:prop:race} is an algorithm as well as a theorem. Given the current state,
compute the rate $\theta_i$ of every possible event and their sum $\Theta$; draw
$T\sim\mathrm{Exp}(\Theta)$ by inverting the distribution function of
\cref{m:def:exp} on a uniform variate; draw the event $k$ with probability
$\theta_k/\Theta$; advance the clock by $T$, apply event $k$, and repeat. Because
the two draws are independent and the exponential law is memoryless, the resulting
trajectory has exactly the law of the process. This is Gillespie's direct method
\cite{Gillespie1977}; it is used in this thesis to check analytical results, never
to replace them.
\end{remark}

\subsection{Continuous-time Markov chains}
\label{m:sec:ctmc}

Let $\{X_t : t\ge0\}$ be a stochastic process taking values in a countable state
space $\mathcal S$, which throughout will be a set of population counts or of
tuples of counts. The process is a \emph{continuous-time Markov chain} if, for all
times $s,t\ge0$ and all states,
\begin{equation}
\label{m:eq:markov}
\pr{X_{t+s} = j \mid X_t = i,\ \{X_u\}_{u<t}} = \pr{X_{t+s}=j\mid X_t = i},
\end{equation}
so that the past influences the future only through the present. The chain is
\emph{time homogeneous} if the right-hand side depends on $s$ but not on $t$, in
which case one writes $p_{ij}(s) = \pr{X_{t+s}=j\mid X_t=i}$. All chains in this
chapter are time homogeneous
\fwd{inhomog}{, and the time-inhomogeneous case is taken up in the later chapter on
non-constant rates}.

Such a chain is specified by its \emph{generator}, the matrix $Q = (q_{ij})$ whose
off-diagonal entries are the rates of the corresponding transitions,
\begin{equation}
\label{m:eq:generator}
q_{ij} = \lim_{\Delta t\to0}\frac{p_{ij}(\Delta t)}{\Delta t}\quad(j\neq i),
\qquad
q_{ii} = -\sum_{j\neq i}q_{ij} =: -q_i ,
\end{equation}
so that every row of $Q$ sums to zero. Equivalently, and this is the form used to
define models below,
\begin{equation}
\label{m:eq:infinitesimal}
\pr{X_{t+\Delta t} = j \mid X_t = i} = q_{ij}\,\Delta t + o(\Delta t)
\qquad (j \neq i),
\end{equation}
with the complementary probability $1-q_i\Delta t + o(\Delta t)$ of no transition
at all. By \cref{m:prop:race} the chain leaves state $i$ after a holding time
distributed as $\mathrm{Exp}(q_i)$ and then jumps to $j$ with probability
$q_{ij}/q_i$. A state with $q_i=0$ is \emph{absorbing}: the chain that reaches it
stays there forever. Extinction and rupture will both be modelled as absorbing
states, and the whole of \cref{m:sec:qsd} concerns the behaviour of a chain
conditioned on not yet having been absorbed.

Differentiating the Chapman--Kolmogorov identity
$p_{ij}(t+\Delta t) = \sum_k p_{ik}(t)p_{kj}(\Delta t)$ with respect to $\Delta t$
at $\Delta t=0$ gives the \emph{forward Kolmogorov equations},
\begin{equation}
\label{m:eq:kolmogorov}
\ddt{p_{ij}(t)} = \sum_{k} p_{ik}(t)\,q_{kj},
\qquad\text{that is}\qquad
\dda{P}{t} = P\,Q .
\end{equation}
Differentiating the same identity in the other slot gives the \emph{backward}
equations $\dd P/\dd t = QP$, which condition on the first jump out of the initial
state rather than on the last jump into the final one; \cref{m:sec:firststep} uses
that form. Written out for a fixed initial state and with $p_j(t)$ for the
probability of occupying state $j$ at time $t$, \cref{m:eq:kolmogorov} becomes the
master equation
\begin{equation}
\label{m:eq:master}
\ddt{p_j(t)} = \sum_{k\neq j}q_{kj}\,p_k(t) - q_j\,p_j(t):
\end{equation}
probability flows into state $j$ from every state that can reach it and out of $j$
at total rate $q_j$. Every model in this thesis is specified by writing down
\cref{m:eq:infinitesimal} and then reading off \cref{m:eq:master}.

\subsection{First-step analysis}
\label{m:sec:firststep}

When the question is not the full time course but a hitting probability or a mean
absorption time, there is a shortcut. Conditioning on the first jump out of the
current state replaces the evolution problem by a linear system, because after that
jump the chain starts afresh from wherever it landed.

\begin{proposition}[First-step principle]
\label{m:prop:firststep}
Let $\{X_t\}$ be a continuous-time Markov chain with absorbing set
$\mathcal A\subset\mathcal S$, and let $h_i = \pr{X_t\in\mathcal A\text{ for some
}t\mid X_0=i}$. Then $h_i=1$ for $i\in\mathcal A$, and for every transient $i$,
\begin{equation}
\label{m:eq:firststep}
h_i = \sum_{j\neq i}\frac{q_{ij}}{q_i}\,h_j .
\end{equation}
\end{proposition}

\begin{proof}
From state $i$ the chain waits an $\mathrm{Exp}(q_i)$ time and then, by
\cref{m:prop:race}, jumps to $j$ with probability $q_{ij}/q_i$, independently of the
waiting time. Absorption occurs eventually if and only if it occurs from the state
reached; conditioning on that state and using the Markov property
\cref{m:eq:markov} gives \cref{m:eq:firststep}.
\end{proof}

Only the jump chain enters \cref{m:eq:firststep}: hitting \emph{probabilities} do
not depend on the holding times, only on which transition wins each race. Mean
hitting times do depend on them, and satisfy the inhomogeneous companion system
$k_i = 1/q_i + \sum_{j\neq i}(q_{ij}/q_i)k_j$, obtained by the same conditioning
with the expected holding time added.

The discrete-time version of the argument is the one used most often below. If the
chain moves in generations, conditioning on the first generation and using
independence of the resulting subpopulations produces a fixed-point equation rather
than a linear system: for the Galton--Watson process of \cref{m:sec:gw} the
extinction probability from a single individual satisfies $s=\phi(s)$, where $\phi$
is the offspring generating function. The three appearances of the idea --- the
clock race, the linear system \cref{m:eq:firststep}, and the branching fixed point
--- are one gesture applied to three state spaces.

\subsection{Probability generating functions}
\label{m:sec:pgf}

For a \rv{} $X$ taking values in $\{0,1,2,\dots\}$ with $p_k = \pr{X=k}$, the
\pgf{} is
\begin{equation}
\label{m:eq:pgfdef}
\phi(z) = \ex{z^X} = \sum_{k=0}^{\infty}p_k z^k ,
\end{equation}
which converges at least for $|z|\le1$ since the $p_k$ sum to one. The series
determines the distribution --- $p_k = \phi^{(k)}(0)/k!$ --- and moments are read
off at $z=1$:
\begin{equation}
\label{m:eq:pgfmoments}
\phi(1) = 1,
\qquad
\phi'(1) = \ex{X},
\qquad
\phi''(1) = \ex{X(X-1)} .
\end{equation}
Two structural properties do most of the work. If $X$ and $Y$ are independent then
$\phi_{X+Y} = \phi_X\phi_Y$, so sums become products. And if $N$ is an independent
count with \pgf{} $\psi$, the random sum $X_1+\dots+X_N$ of independent copies of
$X$ has \pgf{} $\psi\circ\phi$; composition of generating functions is what makes
branching processes tractable, and \cref{m:sec:gw} rests on it.

For a pair of counts the definition extends to
\begin{equation}
\label{m:eq:pgfbivariate}
G(x,y) = \ex{x^{X}y^{Y}} = \sum_{i}\sum_{j}\pr{X=i,\,Y=j}\,x^iy^j ,
\end{equation}
with marginals recovered by setting the other variable to one. The absorption
models of \cref{m:sec:moc} track particles inside and outside a compartment and
therefore need this bivariate form; \cref{m:app:extract} records how to invert it.
The same construction with $n$ arguments,
$G(z_1,\dots,z_n) = \ex{z_1^{X^{(1)}}\cdots z_n^{X^{(n)}}}$, handles vector-valued
states, and is the prepared ground for the multi-type catastrophe models of later
chapters; it is not developed here.

When the counts evolve in time, $G$ acquires a time argument,
$G(x,y,t) = \ex{x^{X_t}y^{Y_t}}$, and the master equation \cref{m:eq:master} for the
infinite family of state probabilities collapses to a single partial differential
equation for $G$. The mechanism is uniform: multiply the equation for
$p_{ij}(t)$ by $x^iy^j$, sum over all states, shift indices so that every sum is
again of the form \cref{m:eq:pgfbivariate}, and recognise the factors of $i$ and $j$
produced by the linear rates as derivatives $\partial_x$ and $\partial_y$. Because
the rates in these models are linear in the counts, the resulting equation is
always of first order, and the method of characteristics therefore always applies
in principle. \Cref{m:sec:linbd} carries this out in the simplest case;
\cref{m:sec:moc} does it three times over.

\subsection{The linear birth--death process}
\label{m:sec:linbd}

The linear birth--death process is the continuous-time counterpart of the
Galton--Watson process, and it is the last piece of standard machinery needed.
Let $X_t$ count individuals, each of which independently gives birth at rate
$\lambda$ and dies at rate $\mu$. In the notation of \cref{m:eq:infinitesimal} the
non-zero rates out of state $n$ are
\begin{equation}
\label{m:eq:bdrates}
q_{n,n+1} = \lambda n,
\qquad
q_{n,n-1} = \mu n ,
\end{equation}
so that $q_n = (\lambda+\mu)n$ and $n=0$ is absorbing. The master equation
\cref{m:eq:master} reads
\begin{equation}
\label{m:eq:bdmaster}
\ddt{p_n(t)} = \lambda(n-1)p_{n-1}(t) + \mu(n+1)p_{n+1}(t) - (\lambda+\mu)n\,p_n(t).
\end{equation}
Multiplying by $z^n$ and summing over $n\ge0$, the three terms become
$\lambda z^2\partial_zG$, $\mu\,\partial_zG$ and $-(\lambda+\mu)z\,\partial_zG$
after the index shifts described above, so that
\begin{equation}
\label{m:eq:bdpde}
\pddt{G} = \bigl(\lambda z^2 - (\lambda+\mu)z + \mu\bigr)\pdd{G}{z}
= (\lambda z-\mu)(z-1)\pdd{G}{z},
\qquad G(z,0) = z^{N},
\end{equation}
for a process started from $N$ individuals. The quadratic coefficient factorises
with roots $z=1$ and $z=\mu/\lambda$; the same factorisation, in a different
variable, reappears in \cref{m:sec:absbirthdeath} and is what makes that calculation
possible. Solving \cref{m:eq:bdpde} along its characteristics gives the classical
result of Kendall \cite{Kendall1948},
\begin{equation}
\label{m:eq:bdsolution}
G(z,t) = \lrb{\frac{\mu(1-z) - (\mu-\lambda z)\ee^{-(\lambda-\mu)t}}
{\lambda(1-z) - (\mu-\lambda z)\ee^{-(\lambda-\mu)t}}}^{\!N},
\qquad \lambda\neq\mu .
\end{equation}
Two consequences are needed later. The mean follows from
\cref{m:eq:pgfmoments} or, more quickly, from the fact that each individual
contributes independently:
\begin{equation}
\label{m:eq:bdmean}
\ex{X_t} = N\ee^{(\lambda-\mu)t}.
\end{equation}
The probability of extinction by time $t$ is $p_0(t) = G(0,t)$, that is
\begin{equation}
\label{m:eq:p0t}
p_0(t) =
\begin{cases}
\displaystyle\lrb{\frac{\mu-\mu\ee^{(\mu-\lambda)t}}{\lambda-\mu\ee^{(\mu-\lambda)t}}}^{\!N},
& \lambda\neq\mu,\\[14pt]
\displaystyle\lrb{\frac{\lambda t}{1+\lambda t}}^{\!N},
& \lambda=\mu,
\end{cases}
\end{equation}
the critical case following by taking $\mu\to\lambda$ in the first line. Letting
$t\to\infty$ recovers the familiar threshold: extinction is certain when
$\lambda\le\mu$, and occurs with probability $(\mu/\lambda)^N$ when
$\lambda>\mu$ --- which is also what \cref{m:eq:firststep} gives, since the embedded
jump chain of \cref{m:eq:bdrates} is a simple random walk on the integers with
up-probability $\lambda/(\lambda+\mu)$. \Cref{m:sec:cts} takes \cref{m:eq:p0t} as its
starting point.
